\section*{Chapitre \ref{ch:b1:non-inertial-frames} --- Référentiels non galiléens~; dynamique terrestre}

\begin{solution}{exo:b1:non-inertial-frames:1}
$\tan\alpha = a/g = 0.255$, $\alpha = 14^\circ$ (le fil part vers
l'arrière)~; $T = m\sqrt{g^2 + a^2} = 1.03\,mg$. Au freinage~: mêmes
angle et tension, le fil partant vers l'avant.
\end{solution}

\begin{solution}{exo:b1:non-inertial-frames:2}
$\Omega = 2\pi/4.0 = \qty{1.57}{rad/s}$~; $\Omega^2r = \qty{7.4}{m/s^2} = 0.75g$,
dirigée vers l'extérieur dans le \omterm{def:b1:point-kinematics:frame}{référentiel} tournant~; c'est le siège
(frottement et maintien) qui fournit la force réelle vers l'intérieur.
\end{solution}

\begin{solution}{exo:b1:non-inertial-frames:3}
$\Omega = 2\pi/86164 = \qty{7.29e-5}{rad/s}$~; $\Omega^2R = \qty{0.034}{m/s^2}$,
soit $0.35\%$ de $g$. Pour s'envoler, il faudrait $\Omega^2R = g$~: $T = 2\pi\sqrt{R/g} =
\qty{84}{min}$ --- dix-sept fois plus vite qu'aujourd'hui.
\end{solution}

\begin{solution}{exo:b1:non-inertial-frames:4}
$a_C = 2\Omega v\sin\lambda = 2 \times 7.29\times10^{-5} \times 27.8 \times 0.707
= \qty{2.9e-3}{m/s^2}$~; $F = \qty{1.1}{kN}$. Vers le nord~: poussé vers
l'est, donc sur le rail de droite (l'oriental)~; vers l'est~: poussé vers
le sud, encore le rail de droite.
\end{solution}

\begin{solution}{exo:b1:non-inertial-frames:5}
$z = \Omega^2r^2/2g$~: \qty{2.0}{mm} à \qty{2}{rad/s}, \qty{5.1}{cm} à
\qty{10}{rad/s}. Un paraboloïde concentre un faisceau parallèle en un
point~: le mercure tournant est un miroir parabolique parfait, bon
marché et qui se polit tout seul (mais que l'on ne peut pas incliner).
\end{solution}

\begin{solution}{exo:b1:non-inertial-frames:6}
L'accélération centrifuge $\Omega^2R\cos\lambda$ est dirigée à l'opposé
de l'axe~; sa composante perpendiculaire à la verticale locale vaut $\Omega^2R\cos\lambda
\sin\lambda$, ce qui incline $\vect g$ de $\delta \approx \Omega^2R\sin\lambda\cos\lambda/g
= 0.034 \times 0.5/9.81 = \qty{1.7e-3}{rad} = 6'$. À l'équateur, $g$ est
réduit de $\Omega^2R = \qty{0.034}{m/s^2}$ par la rotation (l'aplatissement
en ajoute davantage).
\end{solution}

\begin{solution}{exo:b1:non-inertial-frames:7}
$\ddot x_E = 2\Omega gt\cos\lambda$~: $x_E = \tfrac13\Omega g\cos\lambda\,t^3$~; $t =
\sqrt{2h/g} = \qty{4.5}{s}$~: $x_E = \tfrac13 \times 7.29\times10^{-5} \times 9.81
\times 0.707 \times 92 = \qty{1.6}{cm}$. Vers l'est~: le haut du puits va
vers l'est plus vite que le fond (rayon plus grand), et la pierre garde
cet excédent --- ou, dans le \omterm{def:b1:point-kinematics:frame}{référentiel} tournant, la \omterm{thm:b1:non-inertial-frames:dynamics}{force de Coriolis}
sur une vitesse dirigée vers le bas pointe vers l'est.
\end{solution}

\begin{solution}{exo:b1:non-inertial-frames:8}
$T_F = 23.93/\sin\lambda$~: Paris \qty{31.8}{h}~; pôle \qty{23.9}{h}~;
équateur infini (pas de précession). À Paris~: $360^\circ/31.8 = 11.3^\circ$ par heure.
\end{solution}

\begin{solution}{exo:b1:non-inertial-frames:9}
$a_C = 2\Omega v\sin\lambda = \qty{2.1e-3}{m/s^2}$, vers la droite. L'air
qui converge de tous côtés est dévié vers la droite, si bien qu'il
s'enroule vers le centre dans le sens inverse des aiguilles d'une montre
(vu d'en haut)~: le cyclone de l'hémisphère nord. Un évier se vide en
quelques secondes et une tornade se forme en quelques minutes~: sur de
telles durées, la déviation de Coriolis est négligeable devant toute
rotation initiale.
\end{solution}

\begin{solution}{exo:b1:non-inertial-frames:10}
Dans le \omterm{def:b1:point-kinematics:frame}{référentiel} de la voiture, la pesanteur apparente
$\vect g - \vect a$ s'incline vers l'arrière. L'air, plus lourd, «~tombe~»
vers l'arrière~; le ballon, plus léger que l'air qu'il déplace, remonte
la verticale apparente --- il penche donc vers l'\emph{avant}, à
l'inverse d'un fil à plomb.
\end{solution}

\begin{solution}{exo:b1:non-inertial-frames:11}
$\Omega^2R = g$~: $\Omega = \qty{0.31}{rad/s}$, $T = \qty{20}{s}$. $a_C = 2\Omega v
= \qty{0.94}{m/s^2}$, environ $0.1g$~: en marchant dans le sens de la
rotation il se sent plus lourd de $10\%$ (la \omterm{thm:b1:non-inertial-frames:dynamics}{force de Coriolis} pointe
vers l'extérieur), et plus léger à contresens. Pour une balle lâchée~:
dans le \omterm{thm:b1:newton-dynamics:laws}{référentiel galiléen} elle va tout droit tandis que le plancher
se courbe vers elle~; sur la durée de chute $\sqrt{2h/g} = \qty{0.64}{s}$, l'écart
est de l'ordre de $\tfrac13\Omega\,g\,t^3 \cdot 2 \approx 2 \times 0.31 \times
9.81 \times 0.26/3 \approx \qty{0.5}{m}$ --- quelques décimètres, à
contresens de la rotation.
\end{solution}

\begin{solution}{exo:b1:non-inertial-frames:12}
$GM/(d - r)^2 - GM/d^2 \approx 2GMr/d^3$ (au premier ordre en $r/d$). Lune~:
$2 \times 6.67\times10^{-11} \times 7.35\times10^{22} \times 6.37\times10^6/
(3.84\times10^8)^3 = \qty{1.1e-6}{m/s^2}$~; Soleil~: $2 \times 1.33\times10^{20}
\times 6.37\times10^6/(1.50\times10^{11})^3 = \qty{5.0e-7}{m/s^2}$~; rapport
$0.46$~: les marées solaires valent presque la moitié des lunaires, bien
que l'attraction du Soleil soit $180$ fois plus grande --- les marées
varient en $M/d^3$.
\end{solution}

\begin{solution}{pb:b1:non-inertial-frames:1}
\textbf{1.} $T = 2\pi\sqrt{67/9.81} = \qty{16.4}{s}$~; $x_0/\ell = 0.045$
($2.6^\circ$)~: c'est petit.

\textbf{2.} $v_{\max} = x_0\omega_0 = 3.0 \times 0.383 = \qty{1.15}{m/s}$~; hauteur
$x_0^2/2\ell = \qty{6.7}{cm}$.

\textbf{3.} $T = mg + mv_{\max}^2/\ell = mg(1 + 0.002)$~: soit $1.002\,mg$.

\textbf{4.} $\tfrac12\rho C_xSv_{\max}^2 = 0.5 \times 1.2 \times 0.05 \times 1.3 =
\qty{0.04}{N}$, contre $m\omega_0^2x_0 = \qty{12}{N}$~: trois cents fois
moins, et pourtant assez pour arrêter le pendule en quelques heures.

\textbf{5.} Un fil long donne une longue période, donc une masse lente et
une traînée relative faible, et rend visible la précession par
oscillation~; une masse lourde stocke beaucoup d'énergie que la traînée
doit manger --- des heures de balancement.

\textbf{6.} Le poids (vertical) et la tension (dans le plan du fil et de
la verticale, c'est-à-dire dans le plan d'oscillation)~: aucune force n'a
de composante perpendiculaire à ce plan.

\textbf{7.} Le plan est fixe parmi les étoiles~; la glace tourne dessous
vers l'est à $\Omega$, donc l'observateur voit le plan tourner dans le
sens des aiguilles d'une montre et revenir à son point de départ au bout
d'un jour sidéral, \qty{23.93}{h}.

\textbf{8.} Chaque demi-oscillation trace une corde légèrement tournée~:
une rosace~; $86164/8.2 \approx 10500$ demi-oscillations, chacune
$0.034^\circ$ plus loin.

\textbf{9.} $\vect\Omega = \Omega(0, \cos\lambda, \sin\lambda)$.

\textbf{10.} Avec $\vect v' = (\dot x, \dot y, 0)$~:
$\vect\Omega\wedge\vect v' = \Omega(-\sin\lambda\,\dot y,\ \sin\lambda\,\dot x,\ -\cos\lambda\,\dot x)$,
donc $-2m\vect\Omega\wedge\vect v'$ a pour composantes horizontales
$2m\Omega\sin\lambda\,(\dot y, -\dot x)$~: seul $\Omega\sin\lambda$ intervient.

\textbf{11.} $\ddot x = -\omega_0^2x + 2\omega_F\dot y$, $\ddot y = -\omega_0^2y -
2\omega_F\dot x$~; $\omega_F = 7.29\times10^{-5} \times 0.753 = \qty{5.5e-5}{rad/s}
\ll \omega_0 = \qty{0.38}{rad/s}$.

\textbf{12.} $\ddot u = -\omega_0^2u + 2\omega_F(\dot y - j\dot x) = -\omega_0^2u -
2j\omega_F\dot u$.

\textbf{13.} $-\alpha^2 - 2\omega_F\alpha + \omega_0^2 = 0$~: $\alpha = -\omega_F \pm
\sqrt{\omega_F^2 + \omega_0^2} \approx -\omega_F \pm \omega_0$.

\textbf{14.} La parenthèse est une oscillation plane (une droite fixe
passant par l'origine si $A = B$ est réel)~; le facteur
$\eu^{-j\omega_Ft}$ fait tourner tout le motif dans le sens des aiguilles
d'une montre à $\omega_F$~: le plan d'oscillation précesse dans ce sens à
$\Omega\sin\lambda$.

\textbf{15.} $T_F = 2\pi/\omega_F = \qty{1.14e5}{s} = \qty{31.8}{h}$~; $11.3^\circ$
par heure.

\textbf{16.} $\sin\lambda = 0$~: la composante verticale locale de
$\vect\Omega$ s'annule~; la \omterm{thm:b1:non-inertial-frames:dynamics}{force de Coriolis} horizontale est nulle pour
un mouvement horizontal.

\textbf{17.} $2m\omega_Fv_{\max} = 2 \times 28 \times 5.5\times10^{-5} \times 1.15 =
\qty{3.5}{mN}$~; le poids vaut \qty{275}{N} ($10^{-5}$ de celui-ci), la
force de rappel $m\omega_0^2x_0 = \qty{12}{N}$~: trois mille fois moins.

\textbf{18.} $\omega_FT/2 = 5.5\times10^{-5} \times 8.2 = \qty{4.5e-4}{rad}$~:
$x_0 \times 4.5\times10^{-4} = \qty{1.4}{mm}$ par demi-oscillation.

\textbf{19.} Des quilles espacées d'un centimètre sur le cercle tombent
toutes les $\sim 7$ demi-oscillations, soit environ une minute.

\textbf{20.} Une poussée de la main donne une vitesse latérale~: la masse
décrit alors une ellipse, dont la lente précession propre (un effet de
l'\omterm{def:b1:wave-propagation:signal}{amplitude} finie) masque celle de Foucault.

\textbf{21.} Une impulsion radiale ajoute de l'énergie dans le plan
d'oscillation sans ajouter de vitesse transversale, et laisse donc le
plan intact~; une composante latérale entraînerait le pendule dans une
ellipse et simulerait ou masquerait la précession.

\textbf{22.} La vitesse de précession mesure $\Omega\sin\lambda$, la
composante verticale de la rotation terrestre~; $\Omega$ étant connu par
l'astronomie, cette vitesse donne la latitude (et réciproquement).

\textbf{23.} $\sin\lambda = 23.93/48 = 0.499$~: $\lambda = 30^\circ$.

\textbf{24.} Son \omterm{def:b1:angular-momentum:L}{moment cinétique} se conserve, donc son axe reste fixe
parmi les étoiles tandis que la Terre tourne dessous~: c'est le
gyrocompas, qui trouve le nord vrai sans aimant.

\textbf{25.} La \omterm{thm:b1:non-inertial-frames:dynamics}{force de Coriolis} $-2m\vect\Omega\wedge\vect v'$~; sa part
horizontale croît comme $\sin\lambda$~; un tour en $23.93\,\text{h}/\sin\lambda$,
soit \qty{31.8}{h} à Paris.
\end{solution}
